
%\needspace{4\baselineskip}
\color{uniblue}{\subsection{Максимальная токовая защита (МТЗ)}}
\color{black}

\begin{enumerate}[label=\arabic{section}.\arabic{subsection}.\arabic{enumi}, labelsep=4pt, leftmargin=0pt, itemindent=57pt, itemsep=0pt, parsep=5pt] % если что можно крутить itemsep
%%%%%%%%%%%%%%%%%%%%%%%%%%%%%%%%%%%%%%%%%%%%%%%%%%%%%%%%%%%  ОБЩИЕ СВЕДЕНИЯ %%%%%%%%%%%%%%%%%%%%%%%%%%%%%%%%%%%%%%%%%%%%%%%%

\item Общие сведения

\begin{enumerate}[label=\arabic{section}.\arabic{subsection}.\arabic{enumi}.\arabic*, labelsep=4pt, leftmargin=0em, itemindent=65pt, parsep=0pt]

\item
Максимальная токовая защита применяется для защиты от междуфазных КЗ на защищаемом объекте, а также в зоне дальнего резервирования. Функция <<МТЗ>> имеет возможность работы с контролем напряжения. Защищаемыми объектами могут быть вводной выключатель, секционный выключатель и отходящая линия.

\item
Основные характеристики функционального блока <<МТЗ>>:
\begin{itemize}
    \item погрешность по углу $\phi_{\textsl{м.ч.}}$ при номинальных значениях тока и напряжения не более 5$^\circ$.
\end{itemize}

\item
Логическая часть реализует группу пусковых органов с независимыми уставками, предусматривает таймеры выдержек времени на срабатывание, логику пусков и блокировок, коммутацию внутренних пусковых сигналов для формирования выходного сигнала автоматического ускорения.

\item
МТЗ может применяться в качестве:
\begin{itemize}
    \item селективной защиты от перегрузок и междуфазных коротких замыканий на линиях среднего и высокого напряжения;
    \item резервной защиты от перегрузок и междуфазных коротких замыканий в силовых трансформаторах;
    \item резервной защиты сети среднего напряжения в случае установки на стороне НН силовых трансформаторов;
    \item защиты от перегрузок и междуфазных  коротких замыканий специальных видов первичного оборудования (шунтирующих реакторов, конденсаторных батарей и т. д.).
\end{itemize}

\item
В составе функционального блока <<МТЗ>> предусмотрены следующие функции:
\begin{itemize}
\item три независимые направленные/ненаправленные ступени максимальной токовой защиты;
\item комбинированный пусковой орган напряжения;
\item блокировка при броске намагничивающего тока;

\ifboolexpr{bool {isOL}}
{\item блокировка логической защиты шин;}
{\item блокировка логической защиты трансформатора;}

\item орган направления мощности.
\end{itemize}

\item
Функциональный блок <<МТЗ>> может быть полностью выведен из работы оперативно путем активации сигнала <<ОВ МТЗ>>.

\end{enumerate}

%%%%%%%%%%%%%%%%%%%%%%%%%%%%%%%%%%%%%%%%%%%%%%%%%%%%%%%%%  СТУПЕНЬ МТЗ %%%%%%%%%%%%%%%%%%%%%%%%%%%%%%%%%%%%%%%%%%%%%%%%

\item Ступени максимальной токовой защиты (МТЗ 1 ст, МТЗ 2 ст, МТЗ 3 ст)

\begin{enumerate}[label=\arabic{section}.\arabic{subsection}.\arabic{enumi}.\arabic*, labelsep=4pt, leftmargin=0em, itemindent=65pt, parsep=0pt]

\item
Ввод каждой ступени МТЗ в работу на этапе параметрирования устройства осуществляется своим программным переключателем SGF1 <<Ввод\_функции>>. Информация о введенной ступени отображается выходным сигналом для каждой ступени (<<МТЗ 1 ст: Ввод>>, <<МТЗ 2 ст: Ввод>>, <<МТЗ 3 ст: Ввод>>).
\item
Каждая из ступеней <<МТЗ>> может быть выведена из работы оперативно путем активации сигнала оперативного вывода (<<ОВ МТЗ 1 ст>>, <<ОВ МТЗ 2 ст>>, <<ОВ МТЗ 3 ст>>), что характеризуется активным состоянием сигнала <<Оперативный вывод>> (<<МТЗ 1 ст: Оперативный вывод>>, <<МТЗ 2 ст: Оперативный вывод>>, <<МТЗ 3 ст: Оперативный вывод>>) на выходе функции.
\item
При активном состоянии входного сигнала <<МТЗ 1 ст на сигн>>, <<МТЗ 2 ст на сигн>>, <<МТЗ 3 ст на сигн>> действие соответствующей ступени МТЗ переводится на сигнал.
\item
В качестве измерительных органов применяется ИО тока максимального действия, которые контролируют превышение фазными токами заданной величины с коэффициентом возврата равным 0,95. Выбор режима работы ИО (независимая или обратнозависимая характеристика) осуществляется уставкой характеристики срабатывания <<Характеристика>>.
\item
При срабатывании ИО активируются соответствующие выходные сигналы для ступени N МТЗ (N=1,2,3):
\begin{itemize}
    \item <<МТЗ N ст.: ИО IА>>;
    \item <<МТЗ N ст.: ИО IВ>>;
    \item <<МТЗ N ст.: ИО IС>>.
\end{itemize}
\item
При срабатывании ИО и выполнении условий логики работы схемы далее формируются выходные сигнала пуска:
\begin{itemize}
    \item <<МТЗ 1 ст.: Пуск>>;
    \item <<МТЗ 2 ст.: Пуск>>;
    \item <<МТЗ 3 ст.: Пуск>>.
\end{itemize}
\item
При наличии пусковых условий функции запускается таймер Т1 зависимой/независимой выдержки времени на срабатывание. По истечении выдержки времени Т1, в течение которой величина тока постоянно превышала уровень тока срабатывания, активируются соответствующие выходные сигналы срабатывания:
\begin{itemize}
    \item <<МТЗ 1 ст.: Срабатывание>>;
    \item <<МТЗ 2 ст.: Срабатывание>>;
    \item <<МТЗ 3 ст.: Срабатывание>>.
\end{itemize}

\item
При выборе режима с зависимыми характеристиками срабатывания активируется измерительный орган с зависимой выдержкой времени. Все характеристики (кроме характеристики RXIDG) в диапазоне значений тока от $ 1,0 \cdot I/I_{\textsl{ср}} $ до $ 1,1 \cdot I/I_{\textsl{ср}} $ имеют независимую выдержку времени на срабатывание ($t_{\textsl{ср.независ.}}= f (I = 1,1 \cdot I/I_{\textsl{ср}})$). При кратности тока выше 20 ($ I/I_{\textsl{ср}} \geq 20 $)  все зависимые характеристики переходят в независимые (кроме характеристики RXIDG) с выдержкой времени срабатывания, равной выдержки времени срабатывания при $ I/I_{\textsl{ср}} = 20 $ ($t_{\textsl{ср.независ.20}}= f (I = 20 \cdot I/I_{\textsl{ср}})$). C помощью временного коэффициента характеристики срабатывания ($K$) имеется возможность регулирования времени срабатывания времятоковых характеристик.

\item
Графики реализованных в устройстве обратнозависимых времятоковых характеристик приведены в приложении \ref{app:chars}.

\item
Ступень МТЗ имеет возможность выбора направленности. Направленность ступени определяется уставкой SGF3.  Ступень может работать как без контроля направленности, так и с контролем прямого либо обратного направления. Сигналы направления ступень получает от органа направления мощности <<ОНМ>>.
\item
Режим контроля направленности при неисправности цепей напряжения задается уставкой SGF4:
\begin{itemize}
    \item если уставка выставлена в положение <<Блокировка>>,  то при наличии активного сигнала неисправности ЦН осуществляется блокировка всех режимов направленности; 
    \item если уставка выставлена в положение <<Деблокировка>>,  то при наличии активного сигнала неисправности ЦН осуществляется разрешение работы во всех режимах направленности. 
\end{itemize}

\item
Ступень может работать с  пуском по напряжению, который формируется алгоритмом комбинированного пускового органа напряжения <<КПОН>>. 

\item
Порядок работы ступени при неисправности цепей напряжения в режиме работы ступени от КПОН регулируется уставкой SGF8, которая имеет два положения:
\begin{itemize}
    \item деблокировка (чувствительная уставка); 
    \item блокировка (грубая уставка). 
\end{itemize}

\item
Для ступени МТЗ возможны два режима работы с контролем напряжения: <<Вольтметровая блокировка>> и <<Управляющее напряжение>> (выбор режима работы осуществляется программным переключателем SGF6 <<Тип\_КПОН>>). Для выбора режима работы без контроля напряжения необходимо перевести программный переключатель SGF9 <<Реж\_КПОН>> в состояние <<Не предусмотрено>>.
\item
В режиме <<Вольтметровая блокировка>> при превышении уставки по току для срабатывания через выдержку времени таймера T1 (уставка <<Тср>>) необходимо подтверждение от КПОН (снижение напряжения на секции НН). При неисправностях в цепях напряжения ступень МТЗ может или выводиться из работы, или действовать на отключение, в зависимости от выбранного состояния программного переключателя SGF8 <<Реж\_БНН\_КПОН>>: <<Деблокировка>> или <<Блокировка>>. 
\item
Для работы в режиме <<Управляющее напряжение>> необходимо задать две уставки по току для ступени МТЗ: <<Iср>> -- уставка по току срабатывания и <<Iср\_груб>> -- уставка по ток срабатывания грубого органа. В указанном режиме при отсутствии неисправности цепей напряжения при превышении уставки по току <<Iср>> для срабатывания необходимо подтверждение от КПОН (снижение напряжения). При неисправностях в цепях напряжения и при превышении уставки по току грубого органа <<Iср\_груб>> через выдержку времени таймера T1 (уставка T1 <<Тср>>) выдается сигнал срабатывания. Режим работы на уставке <<Iср\_груб>> определяется состоянием программного переключателя SGF8 <<Реж\_БНН\_КПОН>>: <<Деблокировка>> или <<Блокировка>>.

\item
Для ступени МТЗ предусмотрена возможность включения на <<холодную>> нагрузку. Данный режим вводится уставкой SGF2 <<Ввод\_ФВН>> и применяется для временного удвоения уставок по току срабатывания ступени (например, для присоединений с двигательной нагрузкой во время их постановки под нагрузку), для этого в ступени предусмотрен  контроль состояния  дискретного  входа  <<В:  Отключен>>.  При  появлении  сигнала  отключенного  положения выключателя через выдержку времени Т2 активируется выходной сигнал об увеличении уставок по току срабатывания  ступеней  МТЗ  –  <<МТЗ N ст.:  ФВН>>.  При  последующем  включении  выключателя (пропадании  сигнала  <<В:  Отключен>>)  запускается  выдержка  времени  Т3,  в  течение  которой  будет действовать удвоение уставок, после чего произойдет сброс триггера и уставки ступени вернутся к первоначальным значениям. 
\item
Вывод направленности ступени при автоматическом ускорении определяется уставкой SGF5 <<Ненапр\_АУ>>. 
\item
При необходимости возможна блокировка ступени при броске тока намагничивания силового трансформатора от функции <<БНТ>>, это достигается путем перевода программного переключателя выбора режима контроля от БНТ SGF7 <<Реж\_БНТ>> в положение <<Предусмотрено>>.

\item
Алгоритм работы ступеней МТЗ выполнен в соответствии с рисунком \ref{mtz:img1}. В таблице \ref{mtz:tbl1} приведены параметры, необходимые для настройки ступени МТЗ.

\vspace{3mm}
\begin{figure}[H]
\centering
\includegraphics[width=0.95\textwidth,height=0.95\textheight,keepaspectratio]{\fbpath/104. МТЗ ОЛ/_latex/img/mtzst.pdf}
\captionof{figure}{Функциональная схема ступеней <<МТЗ 1 ст>>, <<МТЗ 2 ст>>, <<МТЗ 3 ст>>}\label{mtz:img1}
\end{figure}

\small
\begin{longtable}{|>{\centering\arraybackslash}m{5.3cm}|>{\centering\arraybackslash}m{3.3cm}|>{\centering\arraybackslash}m{4.2cm}|>{\centering\arraybackslash}m{1.8cm}|>{\centering\arraybackslash}m{1cm}|}
\caption{Параметры для настройки ступени МТЗ\hfill\vspace{-0.5\baselineskip}}\label{mtz:tbl1}\\ 
\hline
\rowcolor{gray!30}
Параметр (Параметр на ИЧМ) & Условное обозначение на схеме & Значение/ Диапазон & Единица измерения & Шаг \\ 
\hline
\endfirsthead
\caption*{\hspace{3pt}\emph{Продолжение таблицы \ref{mtz:tbl1}\hfill\vspace{-0.5\baselineskip}}} \\ % сделано по ГОСТ 2.105 п.6.8.7
\hline
\rowcolor{gray!30}
Параметр (Параметр на ИЧМ) & Условное обозначение на схеме & Значение/ Диапазон & Единица измерения & Шаг \\ 
%\hline
\endhead
%\hline
\endfoot
%\hline
\endlastfoot
%==+t1*PTOC1|LVTUVTOC
\centering Ввод функции в работу (Ввод\_функции) & \centering SGF1 & \centering 0 = Не предусмотрено\\1 = Предусмотрено & \centering -- & \centering \arraybackslash -- \\
\hline
\centering Ток срабатывания (Iср) & \centering I> & \centering 0,10 ... 30,00 & \centering о.е. & \centering \arraybackslash 0,01 \\
\hline
\centering Ток срабатывания грубого органа в режиме управляющего напряжения (Iср\_груб.) & \centering Iгруб> & \centering 0,10 ... 30,00 & \centering о.е. & \centering \arraybackslash 0,01 \\
\hline
\centering Ток срабатывания функции включения под нагрузку (Iфвн) & \centering -- & \centering 0,10 ... 30,00 & \centering о.е. & \centering \arraybackslash 0,01 \\
\hline
\centering Ввод функции включения на «холодную» нагрузку (Ввод\_ФВН) & \centering SGF2 & \centering 0 = Не предусмотрено\\1 = Предусмотрено & \centering -- & \centering \arraybackslash -- \\
\hline
\centering Время отключенного состояния выключателя (Тзагруб.) & \centering T2 & \centering 0 ... 20000 & \centering мс & \centering \arraybackslash 10 \\
\hline
\centering Время ввода загрубления (Тсброс.загруб.) & \centering T3 & \centering 0 ... 20000 & \centering мс & \centering \arraybackslash 10 \\
\hline
\centering Тип пуска по напряжению (Тип\_КПОН) & \centering SGF6 & \centering 1 = Управляющее напряжение\\2 = Вольтметровая блокировка & \centering -- & \centering \arraybackslash -- \\
\hline
\centering Режим контроля от БНТ (Реж\_БНТ) & \centering SGF7 & \centering 0 = Не предусмотрено\\1 = Предусмотрено & \centering -- & \centering \arraybackslash -- \\
\hline
\centering Режим направленности (Направленность) & \centering SGF3 & \centering 1 = Ненаправленная\\2 = Прямонаправленная\\3 = Обратнонаправленная & \centering -- & \centering \arraybackslash -- \\
\hline
\centering Режим ОНМ при неисправности ЦН (Реж\_БНН\_ОНМ) & \centering SGF4 & \centering 0 = Блокировка\\1 = Деблокировка & \centering -- & \centering \arraybackslash -- \\
\hline
\centering Вывод направленности при АУ (Ненапр\_АУ) & \centering SGF5 & \centering 0 = Не предусмотрено\\1 = Предусмотрено & \centering -- & \centering \arraybackslash -- \\
\hline
\centering Режим КПОН при неисправности ЦН (Реж\_БНН\_КПОН) & \centering SGF8 & \centering 1 = Деблокировка\\2 = Блокировка & \centering -- & \centering \arraybackslash -- \\
\hline
\centering Режим контроля от КПОН (Реж\_КПОН) & \centering SGF9 & \centering 0 = Не предусмотрено\\1 = Предусмотрено & \centering -- & \centering \arraybackslash -- \\
\hline
\centering Характеристика срабатывания (Характеристика) & \centering -- & \centering 1 = Кривая F\\2 = Кривая E\\4 = Кривая D\\10 = Кривая B\\11 = Кривая A\\12 = Кривая C\\16 = Независимая\\17 = Кривая RXIDG & \centering -- & \centering \arraybackslash -- \\
\hline
\centering Коэффициент регулировки времени срабатывания зависимой характеристики (К) & \centering -- & \centering 0,05 ... 15,00 & \centering -- & \centering \arraybackslash 0,01 \\
\hline
\centering Независимая выдержка времени срабатывания (Tср) & \centering T1 & \centering 0 ... 100000 & \centering мс & \centering \arraybackslash 10 \\
\hline
%===t1
\end{longtable}
\normalsize

\end{enumerate}

%%%%%%%%%%%%%%%%%%%%%%%%%%%%%%%%%%%%%%%%%%%%%%%%%%%%%%%%% ОРГАН ОНМ %%%%%%%%%%%%%%%%%%%%%%%%%%%%%%%%%%%%%%%%%%%%%%%%

\item Орган направления мощности (ОНМ)

\begin{enumerate}[label=\arabic{section}.\arabic{subsection}.\arabic{enumi}.\arabic*, labelsep=4pt, leftmargin=0em, itemindent=65pt, parsep=0pt]

\item
Функция <<ОНМ>> автоматически вводится в работу, если введена в работу одна из ступеней МТЗ и в этой ступени настроен контроль направленности.
\item
Определение направления мощности производится по 90-градусной схеме и использует в качестве воздействующих величин действующие значения фазных токов и действующие значения линейных
напряжений $I_A$ и $U_{BC}$, $I_B$ и $U_{CA}$ и $I_C$ и $U_{AB}$. Название схемы определяется углами сдвига между приложенными к устройству напряжением и током в условиях симметричной трёхфазной системы, где токи каждой фазы соответствуют аналогичным значениям фазных напряжений. Направление устанавливается исходя из величины углового сдвига фазы между током и напряжением индивидуально для каждой пары сигналов.

\item
Направленный элемент срабатывает при выполнении следующих условий:
\begin{itemize}
    \item линейное напряжение превышает 1 вольт;
    \item фазный ток превышает значение $0,04 \cdot I_{\textsl{ном}}$;
    \item вектор полной мощности, определяемый по углу между фазным током и соответствующим ему линейным напряжением, находится в области срабатывания (прямое или обратное направление).
\end{itemize}

\item
Для определения области срабатывания задается уставка $\varphi_{\textsl{м.ч.}}$. Отсчет угла максимальной чувствительности производится от вектора линейного напряжения и считается положительным, если направление отсчета совпадает с направлением вращения против часовой стрелки. Угол сектора отсчитывается от угла максимальной чувствительности в обе стороны и составляет 85°.

\item 
Сигнал <<ОНМ: Прямое>> появляется, если вектор полной мощности находится в диапазоне от $\varphi_{\textsl{м.ч.}}$ - 85° до $\varphi_{\textsl{м.ч.}}$ + 85°. Сигнал <<ОНМ: Обратное>> появляется, если вектор полной мощности находится в диапазоне от ($\varphi_{\textsl{м.ч.}}$ - 180°) + 85° до ($\varphi_{\textsl{м.ч.}}$ - 180°) - 85°. Если вектор полной мощности находится в зоне <<нечувствительности>>, то ни один из вышеуказанных сигналов не появляется.

\item 
Для правильной работы в случае близких коротких замыканий в функции <<ОНМ>> предусмотрено запоминание угла мощности за два периода перед пуском токового измерительного органа на протяжении 40 мс. Запоминание угла мощности осуществляется сразу после поступления пускового сигнала и возможно лишь при условии, что напряжение превышает 1 В до момента запуска. Зафиксированный угол мощности применяется ОНМ только в том случае, если напряжение опускается ниже 0,9 В. В случае активации БНН режим работы <<по памяти>> блокируется. Длительность работы блока по предшествующему значению мощности в этом случае осуществляется в течение 300 мс. Поясняющая диаграмма определения направления мощности приведена на рисунке \ref{mtz:imgonm}.

\vspace{3mm}
\begin{figure}[H]
\centering
\includegraphics[width=0.8\textwidth,height=0.8\textheight,keepaspectratio]{\fbpath/104. МТЗ ОЛ/_latex/img/char.pdf}
\captionof{figure}{Характеристика органа направления мощности}\label{mtz:imgonm}
\end{figure}

\item
Алгоритм работы органа направления мощности (<<ОНМ>>) выполнен в соответствии с рисунком \ref{mtz:imgonm1}. Алгоритм работы логики определения направления, входящий в состав органа направления, выполнен в соответствии с рисунком \ref{mtz:imgonm2}. В таблице \ref{mtz:tblonm} приведены параметры, необходимые для настройки функции <<ОНМ>>.

\small
\begin{longtable}{|>{\centering\arraybackslash}m{5.3cm}|>{\centering\arraybackslash}m{3.3cm}|>{\centering\arraybackslash}m{4.2cm}|>{\centering\arraybackslash}m{1.8cm}|>{\centering\arraybackslash}m{1cm}|}
\caption{Параметры для настройки функции <<ОНМ>>\hfill\vspace{-0.5\baselineskip}}\label{mtz:tblonm}\\ 
\hline
\rowcolor{gray!30}
Параметр (Параметр на ИЧМ) & Условное обозначение на схеме & Значение/ Диапазон & Единица измерения & Шаг \\ 
\hline
\endfirsthead
\caption*{\hspace{3pt}\emph{Продолжение таблицы \ref{mtz:tblonm}\hfill\vspace{-0.5\baselineskip}}} \\ % сделано по ГОСТ 2.105 п.6.8.7
\hline
\rowcolor{gray!30}
Параметр (Параметр на ИЧМ) & Условное обозначение на схеме & Значение/ Диапазон & Единица измерения & Шаг \\ 
%\hline
\endhead
%\hline
\endfoot
%\hline
\endlastfoot
%==+t1*RDIR1|LVTUVTOC
\centering Угол максимальной чувствительности (Фмч) & \centering φмч & \centering -179,0 ... 180,0 & \centering град. & \centering \arraybackslash 0,1 \\
\hline
%===t1
\end{longtable}
\normalsize

\vspace{3mm}
\begin{figure}[H]
\centering
\includegraphics[width=0.8\textwidth,height=0.8\textheight,keepaspectratio]{\fbpath/104. МТЗ ОЛ/_latex/img/onm.pdf}
\captionof{figure}{Функциональная схема <<ОНМ>>}\label{mtz:imgonm1}
\end{figure}

\vspace{3mm}
\begin{figure}[H]
\centering
\includegraphics[width=0.9\textwidth,height=0.9\textheight,keepaspectratio]{\fbpath/104. МТЗ ОЛ/_latex/img/onm1.pdf}
\captionof{figure}{ Функциональная схема логики определения направления}\label{mtz:imgonm2}
\end{figure}

\end{enumerate}

%%%%%%%%%%%%%%%%%%%%%%%%%%%%%%%%%%%%%%%%%%%%%%%%%%%%%%%%% ОРГАН КПОН %%%%%%%%%%%%%%%%%%%%%%%%%%%%%%%%%%%%%%%%%%%%%%%%

\item Комбинированный пусковой орган напряжения (КПОН)

\begin{enumerate}[label=\arabic{section}.\arabic{subsection}.\arabic{enumi}.\arabic*, labelsep=4pt, leftmargin=0em, itemindent=65pt, parsep=0pt]

\item
Функция <<КПОН>> автоматически вводится в работу, если введена в работу одна из ступеней МТЗ и в этой ступени введен режим контроля от КПОН.
\item
Положением программного переключателя SGF1 <<Реж\_пуска>> осуществляется выбор алгоритма работы:
\begin{itemize}
\item по снижению минимального из трех линейных напряжений $U_{ab}, U_{bc}, U_{ca}$ ниже уставки $U_{\textsl{мин}}$ (состояние программного переключателя <<По Uмин>>);
\item комбинированный: по снижению минимального из трех линейных напряжений $U_{ab}, U_{bc}, U_{ca}$ ниже уставки $U_{\textsl{мин}}$ или по повышению напряжения обратной последовательности $U2$ выше уставки $U2_{\textsl{макс}}$ (состояние программного переключателя <<Комбинированный>>);
\item при приеме сигнала от внешнего КПОН <<КПОН внеш.>> (состояние программного переключателя <<Внешний>>).
\end{itemize}
\item
Функция <<КПОН>> контролирует снижение любого из линейных напряжений (\textsl{UAB, UBC, UCA}) ниже заданной уставки или, при введенном положении <<Комбинированный>> программного переключателя <<SGF1>>, повышение напряжения обратной последовательности (\textsl{U2}) выше заданной уставки или прием сигнала <<КПОН внеш.>>, в случае выполнения заданных условий выдается сигнал пуска КПОН.
\item
В качестве воздействующих величин функции <<КПОН>> использует действующие значения основной гармоники линейных напряжений и напряжение обратной последовательности. Коэффициент возврата пороговых элементов равен 1,05 для ИО по минимальному напряжению и 0,95 для ИО по максимальному напряжению.
\item
Алгоритм работы функции <<КПОН>> выполнен в соответствии с рисунком \ref{mtzvn:img2}. В таблице \ref{mtzvn:tbl2} приведены параметры, необходимые для настройки функции <<КПОН>>.

\vspace{3mm}
\begin{figure}[H]
\centering
\includegraphics[width=0.8\textwidth,height=0.8\textheight,keepaspectratio]{\fbpath/104. МТЗ ОЛ/_latex/img/kpon.pdf}
\captionof{figure}{Функциональная схема <<КПОН>>}\label{mtzvn:img2}
\end{figure}

\small
\begin{longtable}{|>{\centering\arraybackslash}m{5.3cm}|>{\centering\arraybackslash}m{3.3cm}|>{\centering\arraybackslash}m{4.2cm}|>{\centering\arraybackslash}m{1.8cm}|>{\centering\arraybackslash}m{1cm}|}
\caption{Параметры для настройки функции <<КПОН>>\hfill\vspace{-0.5\baselineskip}}\label{mtzvn:tbl2}\\ 
\hline
\rowcolor{gray!30}
Параметр (Параметр на ИЧМ) & Условное обозначение на схеме & Значение/ Диапазон & Единица измерения & Шаг \\ 
\hline
\endfirsthead
%\caption{\emph{Продолжение\hfill\vspace{-0.5\baselineskip}}} \\ % Отступ обозначения таблицы от таблицы и выравнивание надписи слева
\caption*{\hspace{3pt}\emph{Продолжение таблицы \ref{mtzvn:tbl2}\hfill\vspace{-0.5\baselineskip}}} \\ % сделано по ГОСТ 2.105 п.6.8.7
\hline
\rowcolor{gray!30}
Параметр (Параметр на ИЧМ) & Условное обозначение на схеме & Значение/ Диапазон & Единица измерения & Шаг \\ 
%\hline
\endhead
%\hline
\endfoot
%\hline
\endlastfoot
%==+t1*PTUV1|LVTUVTOC
\centering Напряжение срабатывания (Uср) & \centering U< & \centering 2,0 ... 100,0 & \centering В & \centering \arraybackslash 0,1 \\
\hline
\centering Напряжение срабатывания ОП (U2ср) & \centering U2> & \centering 2,0 ... 50,0 & \centering В & \centering \arraybackslash 0,1 \\
\hline
\centering Режим пуска (Реж\_пуска) & \centering SGF1 & \centering 0 = По Uмин\\1 = Комбинированный\\2 = Внешний & \centering -- & \centering \arraybackslash -- \\
\hline
%===t1
\end{longtable}
\normalsize

\end{enumerate}
%%%%%%%%%%%%%%%%%%%%%%%%%%%%%%%%%%%%%%%%%%%%%%%%%%%%%%%%% ОРГАН БНТ %%%%%%%%%%%%%%%%%%%%%%%%%%%%%%%%%%%%%%%%%%%%%%%%

\item Блокировка при броске намагничивающего тока (БНТ)

\begin{enumerate}[label=\arabic{section}.\arabic{subsection}.\arabic{enumi}.\arabic*, labelsep=4pt, leftmargin=0em, itemindent=65pt, parsep=0pt]

\item
Функция блокировки при броске намагничивающего тока (<<БНТ>>) позволяет заблокировать защиту при броске тока намагничивания, возникающем при постановке силового трансформатора под напряжение. 
\item
Ввод функции <<БНТ>> осуществляется автоматически при активации режим контроля броска тока намагничивания введенной в работу ступени МТЗ.
\item
Условие пуска выполняется при одновременном наличии следующих условий:
\begin{itemize}
\item отношение действующих значений фазных токов второй гармоники (\textsl{IA 2г, IB 2г, IC 2г}) к действующим значениям фазных токов основной гармоники (\textsl{IA, IB, IC}) превышает уставку <<Ih2/Ih1>> (контролируется для каждой фазы отдельно);
\item отсутствует срабатывание порогового элемента, реагирующего на превышение максимального из действующих значений основной гармоники фазных токов (\textsl{IA, IB, IC}) над уставкой <<I\_ср>>;
\item сформирован пуск одной из ступеней МТЗ (контроль выполняется пофазный).
\end{itemize}

\item
В функции предусмотрена возможность ввода перекрестной блокировки: при переводе программного переключателя SGF1 <<Перекрест\_блок>> в состояние <<Предусмотрено>> обнаружение броска тока намагничивания в одной из фаз блокирует срабатывание по всем фазам ступени МТЗ. Пуск функции <<БНТ>> учитывается отдельно для каждой из ступеней МТЗ при помощи соответствующего программного переключателя <<Реж\_БНТ>> в составе ступени МТЗ. При этом блокировка каждой ступени осуществляется пофазно.
\item
Алгоритм работы функции <<БНТ>> выполнен в соответствии с рисунком \ref{mtzvn:img3}. В таблице \ref{mtzvn:tbl3} приведены параметры, необходимые для настройки функции <<БНТ>>.

\vspace{3mm}
\begin{figure}[H]
\centering
\includegraphics[width=0.9\textwidth,height=0.9\textheight,keepaspectratio]{\fbpath/104. МТЗ ОЛ/_latex/img/bnt.pdf}
\captionof{figure}{Функциональная схема <<БНТ>>}\label{mtzvn:img3}
\end{figure}

\small
\begin{longtable}{|>{\centering\arraybackslash}m{5.3cm}|>{\centering\arraybackslash}m{3.3cm}|>{\centering\arraybackslash}m{4.2cm}|>{\centering\arraybackslash}m{1.8cm}|>{\centering\arraybackslash}m{1cm}|}
\caption{Параметры для настройки функции <<БНТ>>\hfill\vspace{-0.5\baselineskip}}\label{mtzvn:tbl3}\\ 
\hline
\rowcolor{gray!30}
Параметр (Параметр на ИЧМ) & Условное обозначение на схеме & Значение/ Диапазон & Единица измерения & Шаг \\ 
\hline
\endfirsthead
\caption*{\hspace{3pt}\emph{Продолжение таблицы \ref{mtzvn:tbl3}\hfill\vspace{-0.5\baselineskip}}} \\ % сделано по ГОСТ 2.105 п.6.8.7
\hline
\rowcolor{gray!30}
Параметр (Параметр на ИЧМ) & Условное обозначение на схеме & Значение/ Диапазон & Единица измерения & Шаг \\ 
%\hline
\endhead
%\hline
\endfoot
%\hline
\endlastfoot
%==+t1*PHAR1|LVTUVTOC
\centering Отношение тока второй гармоники к току основной гармоники (Ih2/Ih1) & \centering I2г/I1г> & \centering 5 ... 100 & \centering \% & \centering \arraybackslash 1 \\
\hline
\centering Ток блокировки (Iср) & \centering I> & \centering 0,1 ... 15,0 & \centering о.е. & \centering \arraybackslash 0,1 \\
\hline
\centering Перекрестная блокировка (Перекрест\_блок) & \centering SGF1 & \centering 0 = Не предусмотрено\\1 = Предусмотрено & \centering -- & \centering \arraybackslash -- \\
\hline
%===t1
\end{longtable}
\normalsize

\end{enumerate}

%%%%%%%%%%%%%%%%%%%%%%%%%%%%%%%%%%%%%%%%%%%%%%%%%%%%%%%%% ОРГАН БЛЗШ %%%%%%%%%%%%%%%%%%%%%%%%%%%%%%%%%%%%%%%%%%%%%%%%
%%% ИСПОЛЬЗУЕМ ЛОГИЧЕСКИЕ УСЛОВИЯ!!!
\ifboolexpr{bool {isOL}}
{
\item Блокировка логической защиты шин (БЛЗШ)
\begin{enumerate}[label=\arabic{section}.\arabic{subsection}.\arabic{enumi}.\arabic*, labelsep=4pt, leftmargin=0em, itemindent=65pt, parsep=0pt]

\item
Функция блокировки логической защиты шин (<<БЛЗШ>>) предусматривается для блокировки ЛЗШ на вводном и секционном выключателях секции. В качестве воздействующих величин используются сигналы пуска ступеней МТЗ. Состояние программного переключателя SGF1 <<БлокЛЗ\_выбор\_ст>> определяет ступень МТЗ, от которой будет производиться блокировка ЛЗШ.
\item
Алгоритм работы <<БЛЗШ>> выполнен в соответствии с рисунком \ref{mtzvn:img4}. В таблице \ref{mtzvn:tbl4} приведены параметры, необходимые для настройки функции <<БЛЗШ>>.

\begin{figure}[H]
\centering
\includegraphics[width=0.6\textwidth,height=0.6\textheight,keepaspectratio]{\fbpath/104. МТЗ ОЛ/_latex/img/blzsh.pdf}
%\includegraphics[scale=0.9,keepaspectratio]{\fbpath/104. МТЗ ОЛ/_latex/img/blzsh.pdf}
\captionof{figure}{Функциональная схема БЛЗШ}\label{mtzvn:img4}
\vspace{3mm}
\end{figure}

\small
\begin{longtable}{|>{\centering\arraybackslash}m{5.3cm}|>{\centering\arraybackslash}m{3.3cm}|>{\centering\arraybackslash}m{4.2cm}|>{\centering\arraybackslash}m{1.8cm}|>{\centering\arraybackslash}m{1cm}|}
\caption{Параметры для настройки функции <<БЛЗШ>>\hfill\vspace{-0.5\baselineskip}}\label{mtzvn:tbl4}\\ 
\hline
\rowcolor{gray!30}
Параметр (Параметр на ИЧМ) & Условное обозначение на схеме & Значение/ Диапазон & Единица измерения & Шаг \\ 
\hline
\endfirsthead
\caption*{\hspace{3pt}\emph{Продолжение таблицы \ref{mtzvn:tbl4}\hfill\vspace{-0.5\baselineskip}}} \\ % сделано по ГОСТ 2.105 п.6.8.7
\hline
\rowcolor{gray!30}
Параметр (Параметр на ИЧМ) & Условное обозначение на схеме & Значение/ Диапазон & Единица измерения & Шаг \\ 
%\hline
\endhead
%\hline
\endfoot
%\hline
\endlastfoot
%==+t1*RBLC1|LVTUVTOC
\centering Выбор ступени блокировки (БлокЛЗ\_выбор\_ст) & \centering SGF1 & \centering 1 = Не предусмотрено\\2 = 1 ступень\\3 = 2 ступень\\4 = 3 ступень & \centering -- & \centering \arraybackslash -- \\
\hline
%===t1
\end{longtable}
\normalsize    
}{}
\ifboolexpr{bool {isVV} or bool {isSV}}{
\item Блокировка логической защиты трансформатора (БЛЗТ)
\begin{enumerate}[label=\arabic{section}.\arabic{subsection}.\arabic{enumi}.\arabic*, labelsep=4pt, leftmargin=0em, itemindent=65pt, parsep=0pt]

\item
Функция блокировки логической защиты трансформатора (<<БЛЗТ>>) предусматривается для блокировки ЛЗТ вышестоящего силового трансформатора. В качестве воздействующих величин используются сигналы пуска ступеней МТЗ. Состояние программного переключателя SGF1 <<БлокЛЗ\_выбор\_ст>> определяет ступень МТЗ, от которой будет производиться блокировка ЛЗТ.
\item
Алгоритм работы <<БЛЗТ>> выполнен в соответствии с рисунком \ref{mtzvn:img4}. В таблице \ref{mtzvn:tbl4} приведены параметры, необходимые для настройки функции <<БЛЗТ>>.

\begin{figure}[H]
\centering
\includegraphics[width=0.6\textwidth,height=0.6\textheight,keepaspectratio]{\fbpath/104. МТЗ ОЛ/_latex/img/blzt.pdf}
\captionof{figure}{Функциональная схема БЛЗТ}\label{mtzvn:img4}
\vspace{3mm}
\end{figure}

\small
\begin{longtable}{|>{\centering\arraybackslash}m{5.3cm}|>{\centering\arraybackslash}m{3.3cm}|>{\centering\arraybackslash}m{4.2cm}|>{\centering\arraybackslash}m{1.8cm}|>{\centering\arraybackslash}m{1cm}|}
\caption{Параметры для настройки функции <<БЛЗТ>>\hfill\vspace{-0.5\baselineskip}}\label{mtzvn:tbl4}\\ 
\hline
\rowcolor{gray!30}
Параметр (Параметр на ИЧМ) & Условное обозначение на схеме & Значение/ Диапазон & Единица измерения & Шаг \\ 
\hline
\endfirsthead
\caption*{\hspace{3pt}\emph{Продолжение таблицы \ref{mtzvn:tbl4}\hfill\vspace{-0.5\baselineskip}}} \\ % сделано по ГОСТ 2.105 п.6.8.7
\hline
\rowcolor{gray!30}
Параметр (Параметр на ИЧМ) & Условное обозначение на схеме & Значение/ Диапазон & Единица измерения & Шаг \\ 
%\hline
\endhead
%\hline
\endfoot
%\hline
\endlastfoot
%==+t1*RBLC1|VV_LVTUVTOC
\centering Выбор ступени блокировки (БлокЛЗ\_выбор\_ст) & \centering SGF1 & \centering 1 = Не предусмотрено\\2 = 1 ступень\\3 = 2 ступень\\4 = 3 ступень & \centering -- & \centering \arraybackslash -- \\
\hline
%===t1
\end{longtable}
\normalsize     
}{}

\end{enumerate}
%%%%%%%%%%%%%%%%%%%%%%%%%%%%%%%%%%%%%%%%%%%%%%%%%%%%%%%%% КОНЕЦ ОРГАН БЛЗШ %%%%%%%%%%%%%%%%%%%%%%%%%%%%%%%%%%%%%%%%%%%%%%%%

%%%%%%%%%%%%%%%%%%%%%%%%%%%%%%%%%%%%%%%%%%%%%%%%%%%%%%%%% ОРГАН АУ %%%%%%%%%%%%%%%%%%%%%%%%%%%%%%%%%%%%%%%%%%%%%%%%

\item Автоматическое ускорение (АУ МТЗ)

\begin{enumerate}[label=\arabic{section}.\arabic{subsection}.\arabic{enumi}.\arabic*, labelsep=4pt, leftmargin=0em, itemindent=65pt, parsep=0pt]

\item
В устройстве предусмотрен режим автоматического ускорения ступеней МТЗ. Выбор ускоряемых ступеней защиты осуществляется из соображений надежного охвата, защищаемого энергообъекта.
\item
Предусмотрена возможность автоматического ускорения для любой ступени МТЗ. Ускоряемая ступень выбирается уставкой SGF1. Автоматическое ускорение срабатывает без выдержки времени.

\item
Алгоритм работы <<АУ МТЗ>> выполнен в соответствии с рисунком \ref{mtz:imgau}. В таблице \ref{mtz:tblau} приведены параметры, необходимые для настройки функции <<АУ МТЗ>>.

\begin{figure}[H]
\centering

\includegraphics[width=0.6\textwidth,height=0.6\textheight,keepaspectratio]{\fbpath/104. МТЗ ОЛ/_latex/img/au.pdf}
\captionof{figure}{Функциональная схема АУ МТЗ}\label{mtz:imgau}
\vspace{3mm}
\end{figure}

\small
\begin{longtable}{|>{\centering\arraybackslash}m{5.3cm}|>{\centering\arraybackslash}m{3.3cm}|>{\centering\arraybackslash}m{4.2cm}|>{\centering\arraybackslash}m{1.8cm}|>{\centering\arraybackslash}m{1cm}|}
\caption{Параметры для настройки функции <<АУ МТЗ>>\hfill\vspace{-0.5\baselineskip}}\label{mtz:tblau}\\ 
\hline
\rowcolor{gray!30}
Параметр (Параметр на ИЧМ) & Условное обозначение на схеме & Значение/ Диапазон & Единица измерения & Шаг \\ 
\hline
\endfirsthead
\caption*{\hspace{3pt}\emph{Продолжение таблицы \ref{mtz:tblau}\hfill\vspace{-0.5\baselineskip}}} \\ % сделано по ГОСТ 2.105 п.6.8.7
\hline
\rowcolor{gray!30}
Параметр (Параметр на ИЧМ) & Условное обозначение на схеме & Значение/ Диапазон & Единица измерения & Шаг \\ 
%\hline
\endhead
%\hline
\endfoot
%\hline
\endlastfoot
%==+t1*AUAPSSU1|LVTUVTOC
\centering Выбор ускоряемой ступени (Выбор\_ст) & \centering SGF1 & \centering 1 = 1 ступень\\2 = 2 ступень\\3 = 3 ступень & \centering -- & \centering \arraybackslash -- \\
\hline
%===t1
\end{longtable}
\normalsize

\end{enumerate}
%%%%%%%%%%%%%%%%%%%%%%%%%%%%%%%%%%%%%%%%%%%%%%%%%%%%%%%%% КОНЕЦ ОРГАН АУ %%%%%%%%%%%%%%%%%%%%%%%%%%%%%%%%%%%%%%%%%%%%%%%%

%%%%%%%%%%%%%%%%%%%%%%%%%%%%%%%%%%%%%%%%%%%%%%%%%%%%%%%%% ОРГАН ОУ %%%%%%%%%%%%%%%%%%%%%%%%%%%%%%%%%%%%%%%%%%%%%%%%

\item Оперативное ускорение (ОУ МТЗ)

\begin{enumerate}[label=\arabic{section}.\arabic{subsection}.\arabic{enumi}.\arabic*, labelsep=4pt, leftmargin=0em, itemindent=65pt, parsep=0pt]

\item
Оперативное ускорение предусматривает ускорение второй или третьей ступени МТЗ. Оперативное ускорение может вводиться:
\begin{itemize}
    \item при выводе основной защиты оборудования;
    \item при опробовании оборудования рабочим напряжением;
    \item при операциях с разъединителями на противоположном конце ВЛ.
\end{itemize}

\item
В данном алгоритме при введенном ключе <<Ввод ОУ МТЗ>> предусмотрена возможность оперативного ускорения второй или третьей ступени МТЗ. Выбор ускоряемой ступени осуществляется программным переключателем SGF1 <<Выбор\_ст\_ОУ>>. При введенной функции и появлению сигнала пуска выбранной ступени по истечению выдержки времени, определяемой уставкой T1, формируется сигнал срабатывания.

\item
Алгоритм работы <<ОУ МТЗ>> выполнен в соответствии с рисунком \ref{mtz:imgou}. В таблице \ref{mtz:tblou} приведены параметры, необходимые для настройки функции <<ОУ МТЗ>>.

\begin{figure}[H]
\centering

\includegraphics[width=0.75\textwidth,height=0.75\textheight,keepaspectratio]{\fbpath/104. МТЗ ОЛ/_latex/img/ou.pdf}
\captionof{figure}{Функциональная схема ОУ МТЗ}\label{mtz:imgou}
\vspace{3mm}
\end{figure}

\small
\begin{longtable}{|>{\centering\arraybackslash}m{5.3cm}|>{\centering\arraybackslash}m{3.3cm}|>{\centering\arraybackslash}m{4.2cm}|>{\centering\arraybackslash}m{1.8cm}|>{\centering\arraybackslash}m{1cm}|}
\caption{Параметры для настройки функции <<ОУ МТЗ>>\hfill\vspace{-0.5\baselineskip}}\label{mtz:tblou}\\ 
\hline
\rowcolor{gray!30}
Параметр (Параметр на ИЧМ) & Условное обозначение на схеме & Значение/ Диапазон & Единица измерения & Шаг \\ 
\hline
\endfirsthead
\caption*{\hspace{3pt}\emph{Продолжение таблицы \ref{mtz:tblou}\hfill\vspace{-0.5\baselineskip}}} \\ % сделано по ГОСТ 2.105 п.6.8.7
\hline
\rowcolor{gray!30}
Параметр (Параметр на ИЧМ) & Условное обозначение на схеме & Значение/ Диапазон & Единица измерения & Шаг \\ 
%\hline
\endhead
%\hline
\endfoot
%\hline
\endlastfoot
%==+t1*OPAPSSU1|LVTUVTOC
\centering Выбор оперативно ускоряемой ступени (Выбор\_ст\_ОУ) & \centering SGF1 & \centering 2 = 2 ступень\\3 = 3 ступень & \centering -- & \centering \arraybackslash -- \\
\hline
\centering Выдержка времени срабатывания оперативного ускорения (T\_ОУ) & \centering T1 & \centering 0 ... 3000 & \centering мс & \centering \arraybackslash 10 \\
\hline
%===t1
\end{longtable}
\normalsize

\end{enumerate}
%%%%%%%%%%%%%%%%%%%%%%%%%%%%%%%%%%%%%%%%%%%%%%%%%%%%%%%%% КОНЕЦ ОРГАН ОУ %%%%%%%%%%%%%%%%%%%%%%%%%%%%%%%%%%%%%%%%%%%%%%%%

%%%%%%%%%%%%%%%%%%%%%%%%%%%%%%%%%%%%%%%%%%%%%%%%%%%%%%%%% ВНЕШНЯЯ ЛОГИКА %%%%%%%%%%%%%%%%%%%%%%%%%%%%%%%%%%%%%%%%%%%%%%%%

\item Внешняя логика функции МТЗ

\begin{enumerate}[label=\arabic{section}.\arabic{subsection}.\arabic{enumi}.\arabic*, labelsep=4pt, leftmargin=0em, itemindent=65pt, parsep=0pt]

\item
Для функции предусмотрен общий выходной сигнал <<Пуск МТЗ>>, который формируется при наличии пуска любой ступени МТЗ. Общий сигнал пуска МТЗ участвует в логике таких функций как: УРОВ, БНН, АУ.

\item
Алгоритм работы внешней логики МТЗ выполнен в соответствии с рисунком \ref{mtz:imgvl}. 

\begin{figure}[H]
\centering

\includegraphics[width=0.65\textwidth,height=0.65\textheight,keepaspectratio]{\fbpath/104. МТЗ ОЛ/_latex/img/vl.pdf}
\captionof{figure}{Функциональная схема внешней логики МТЗ}\label{mtz:imgvl}
\vspace{3mm}
\end{figure}

\end{enumerate}
%%%%%%%%%%%%%%%%%%%%%%%%%%%%%%%%%%%%%%%%%%%%%%%%%%%%%%%%% КОНЕЦ ВНЕШНЯЯ ЛОГИКА %%%%%%%%%%%%%%%%%%%%%%%%%%%%%%%%%%%%%%%%%%%%%%%%

\end{enumerate}
